The most successful model of physics is the Standard Model (SM) of particle physics \cite{Glashow:1961tr, Weinberg:1967tq}. Its description and prediction of elementary particle interactions have been found to be experimentally verified with up to 12 digits of precision \cite{Fan:2022eto}. The predictive power of the SM establishes it as our foundational theory of nature \cite{SLAC-SP-017:1974ind,E598:1974sol,PLUTO:1978jrw, CDF:1995wbb, D0:1995jca,UA1:1983mne, UA1:1983crd, UA2:1983mlz, UA2:1983tsx}. 

Historically, one of the most significant theoretical issues, that the model had to overcome was explaining the generation of mass of force carriers and fermions. A solution to this problem was proposed already in the 1960's by Higgs \cite{Higgs:1964ia, Higgs:1966ev}, which was discovered half a century later \cite{ATLAS:2012yve}, adding to the already successful model. Despite the model's ability to predict physical observables to high accuracy, the model remains an incomplete theory. On the phenomenological side alone, it leaves several fundamental questions unanswered, including: \textit{Why there is a flavour mass hierarchy? Where the matter-antimatter asymmetry comes from? How does Einstein's theory of gravity fit in? And what is dark matter?} 

Furthermore, the model is an effective field theory, meaning the model cannot be valid across all energy scales. This attribute of the model originates from the requirement of having a renormalisable theory as our foundational model. Renormalisation introduces a \textit{running} of the couplings, which greatly affect the Standard Model. The theory of strong interactions (QCD) has a gauge coupling which converges to zero at high energies, however at low energies the coupling diverges, and the famous characteristic of QCD emerges: confinement \cite{Gross:1973id}. This renders perturbative predictions invalid in the IR. In contrast, the electromagnetic interactions (QED) behaves oppositely. Its gauge coupling diverges in the UV leaving us with a incomplete description at high energies. Thus, the Standard Model has neither a free nor an interacting ultraviolet fixed point for all its couplings, leaving us with an effective field theory. 

To address some of these phenomenological questions and extend the SM to higher energies, Grand Unified Theories (GUT)s \cite{Georgi:1974yf, Fritzsch:1974nn, Pati:1974yy, Bai:2021tgl, Langacker:1980js} have been prosed as extensions of the SM. These aim to unify the electromagnetic, weak, and strong force into a single coupling described by a single gauge group. To ensure the validity of these models to energies up to the Planck scale at $10^{19}\,$GeV, where quantum gravitational effects are expected to emerge. The unified gauge coupling in such theories is often \textit{asymptotically free}, while most often the Yukawa and quartic couplings are left untamed in the UV. For gauge-fermion theories without scalars a landscape approach to GUT-model building can be found in \cite{Cacciapaglia:2025nff}.

Theories which are asymptotically free in all couplings, and where the Yukawa and quartic couplings tends to zero at the same rate or faster than the gauge coupling, are known as \textit{complete asymptotically free} (CAF) models \cite{Cheng:1973nv,Gross:1973ju,Holdom_2015,Giudice:2014tma,Pica:2016krb,Callaway:1988ya,Hansen:2017pwe,Zimmermann:1984sx}. These models are valid at arbitrarily high energy, since they have UV fixed points. 

The SM and its extensions are chiral-gauge Yukawa theories, it is therefore paramount to understand UV behaviour of such models. To maintain a general approach, without restricting our focus to one specific model as done in \cite{Giudice:2014tma}, we investigate and construct CAF chiral-gauge Yukawa theories by addressing the following question:

\begin{mdframed}[style=thesisbox]
    \begin{center}
        {\large\textsc{How can we construct chiral-gauge Yukawa theories with different scalar representations such that they remain completely asymptotically free?}}
    \end{center}
\end{mdframed}

In this thesis we will approach this question by considering the generalised Georgi--Glashow (GG)\cite{Georgi:1974sy,Georgi:1974yf} and the Bars--Yankielowicz (BY) \cite{Bars:1981se} models. The IR behaviour of these models in the absences of scalar have been widely investigated in \cite{Appelquist:2000qg,Bolognesi:2020mpe,Csaki:2021xhi,Bolognesi:2021yni,Bolognesi:2021hmg,Bai:2021tgl,Bolognesi:2023xxv,Li:2025tvu,Li:2026ayh}. In this thesis, we will instead consider the high energy behaviour as in \cite{Molgaard:2016bqf}, but extending this work by including a scalar in the adjoint representation. 

We systematically evaluate the impact of adding a single scalar in the fundamental representation and the adjoint representation separately. These two representations of scalars are vital to GUTs. The fundamental scalar plays the role of the Higgs in the SM, breaking the electroweak symmetry and generating massive force carriers and fermions. The adjoint scalar allows us to, at low energies, realise the SM, effectively breaking the gauge group (or a subgroup) of the GUT down to the gauge sector of the SM \cite{Georgi:1974yf, Li:1973mq}. To make the analysis more general, we also consider the changing the number of Chiral families, and the effects of this on the CAF regions. 

This thesis is organised as follows. In Chapter \ref{chap:RG}, we introduce the renormalisation group, and through this beta functions and asymptotic freedom, to lay the ground work of the CAF framework used in this project. In Chapter \ref{chap:GUT}, we further our understanding of the SM, its shortcomings and through this the motivation of GUTs. We will also investigate the $\SU(5)$ GUT to understand the dynamics and principles of GUTs. With this knowledge we can then understand the approach to finding CAF theories described in Chapter \ref{chap:CAF}. Equipped with the CAF framework, we introduce the general Chiral models in Chapter \ref{chap:chiralmodels} and systematically investigate possible CAF regions of the theory. We first consider the models without a scalar in Section \ref{sec:noscalar}, then adding one fundamental scalar in Section \ref{sec:fundscalar}, and finally considering the adjoint scalar in Section \ref{sec:adj}. This analysis leads us to the answer of our question, which is given in Chapter \ref{chap:conc} and includes a summary table stating CAF viability of all considered models.
